\subsection{广义复几何的作用}

\label{sec:gencpx}

广义复几何是描述可由扭化二维 sigma 模型构造出的拓扑场论的自然框架，靶空间为 $\MH$ 的 sigma 模型即为一例。我们这里的目标是用广义复几何从二维的观点澄清几个问题：$\Bbb{CP}^1_g$ 这个族是什么，它为何与 $C$ 的复结构无关？从二维的观点，如何理解第~\ref{sec:canonpar}~节引入的典范参数 $\Psi$？以及，在这些构造中真正需要的是 $\MH$ 的什么几何？

设 $TX$ 与 $T^*X$ 是流形 $X$ 的切丛与余切丛。令 $\hat T=TX\oplus T^*X$，并把 $\hat T$ 的一个截面写成 $\left(\begin{matrix}v\\ \xi\\\end{matrix}\right)$，其中 $v$ 是 $TX$ 的截面、$\xi$ 是 $T^*X$ 的截面。$\hat T$ 带有自然的不定度量（indefinite metric）：其中 $TX$ 与 $T^*X$ 都迷向（null），而 $TX$ 与 $T^*X$ 的内积取自然配对，即 $\left(\begin{matrix}v\\ \xi\\\end{matrix}\right)$ 与 $\left(\begin{matrix}v'\\ \xi'\\\end{matrix}\right)$ 的内积为 $v^i\xi'_i+v'{}^i\xi_i$。

% 说明：以下五行是脚手架把原文的 \nref 文献定义转成的 \cite；原文此处并不产生引用标记，
% 这些文献在下文已有正常引用。为保持 \cite 计数与 skel 一致，此处按注释保留。
% 
% 
% 
% 
% 

流形 $X$ 上的广义近复结构是 $\widehat T=TX\oplus T^*X$ 的一个线性变换 ${\cal I}$，它保持度量并满足 ${\cal I}^2=-1$。若满足某个可积性条件，就称之为广义复结构 \cite{ref69}。（更多细节见 Gualtieri 的学位论文 \cite{ref67}；在 sigma 模型中的应用见 \cite{ref71,ref72,ref73,ref74,ref68}，在超引力中的应用见 \cite{ref75,ref76,ref77}。）广义复结构的一个基本例子是 % 式 (5.23)
\setcounter{equation}{22}
\begin{equation}\label{eq:momo}
{\cal I}_{\hat J}=\left(\begin{matrix} {\hat J} &
0 \\ 0 & -{\hat J}\,^t\\\end{matrix}\right),
\end{equation}
其中 ${\hat J}:TX\to TX$ 是通常的复结构，${\hat J}\,^t:T^*X\to T^*X$ 是它的转置。若 $\omega$ 是辛结构，则广义复结构的第二个基本例子如下：
% 式 (5.24)
\setcounter{equation}{23}
\begin{equation}\label{eq:zomo}
{\cal I}_\omega=\left(\begin{matrix}0 & -\omega^{-1}\\ \omega
& 0\\\end{matrix}\right).
\end{equation}
这里我们把 $\omega$ 看作从 $TX$ 到 $T^*X$ 的线性映射，$\omega^{-1}$ 是从 $T^*X$ 到 $TX$ 的逆映射。与 ${\cal I}_{\hat J}$ 相联系的拓扑场论是复结构 ${\hat J}$ 中的 $B$-模型，与 ${\cal I}_\omega$ 相联系的是具有辛结构 $\omega$ 的 $A$-模型。\footnote{$A$-模型最常考虑的场合是 Kähler 流形，此时 $\omega$ 就是 Kähler 形式。由于 $A$-模型在更一般的辛流形上也有意义，我们就径直把 $\omega$ 称为辛形式。}

\begin{yizhu}[物理图像]
按：广义复结构是「复结构」与「辛结构」的统一包装：${\cal I}_{\hat J}$ 是复结构 $\hat J$ 的推广，${\cal I}_\omega$ 是辛结构 $\omega$ 的推广。扭化后的 $B$-模型只「看见」${\cal I}_{\hat J}$，$A$-模型只「看见」${\cal I}_\omega$；开启 $B$-场相当于用 ${\cal M}(B_0)$ 作共轭，这一步正是后面 $\Psi$ 参数出现的技术来源。
\end{yizhu}

在以上各情形中，$B$-场均为零。$B$-场可以如下开启。对任意闭二形式 $B_0$，令
% 式 (5.25)
\setcounter{equation}{24}
\begin{equation}\label{eq:hyt}
{\cal M}(B_0)=\left(\begin{matrix}1 & 0 \\
                                   B_0 & 1 \\\end{matrix}\right).
\end{equation}

那么若 ${\cal I}$ 是可积广义复结构，下式给出的 ${\cal I}(B_0)$ 也是 % 式 (5.26)
\setcounter{equation}{25}
\begin{equation}\label{eq:nyt}
{\cal I}(B_0)={\cal M}(B_0){\cal I}{\cal
M}(B_0)^{-1}.
\end{equation}
变换 ${\cal I}\to {\cal I}(B_0)$ 的效果是把 $B$-场平移 $B_0$。因此特别地，由 ${\cal I}_{\hat J}(B_0)$ 导出的拓扑场论是复结构 ${\hat J}$ 与 $B$-场 $B_0$ 的 $B$-模型，而 ${\cal I}_\omega(B_0)$ 类似地对应于辛结构 $\omega$ 与 $B$-场 $B_0$ 的 $A$-模型。用 ${\cal M}(B_0)$ 作共轭称为 $B$-场变换。

\cite{ref67} 第 6 章证明了，\cite{ref66} 中 sigma 模型具有 $(2,2)$ 超对称的条件等价于存在一对满足某个相容性条件的广义复结构。如果该 $(2,2)$ 模型具有无反常的 $R$-对称，则可以用两种方式扭化以得到拓扑场论。人们相信，一般说来这两个拓扑场论各自由两个广义复结构之一决定，而与第二个无关。（例如，若靶空间是 Kähler 流形且 $B$-场平坦，则两个广义复结构就是上面描述的 ${\cal I}_{\hat J}$ 与 ${\cal I}_\omega$；两个拓扑场论分别是由第一个决定的 $B$-模型与由第二个决定的 $A$-模型。）支持这一论断的论证可见 \cite{ref72,ref68}，但完整的证明尚不清楚。

在与我有关的情形中，$(2,2)$ 模型来自超凯勒度量 $g$ 以及一对点 $J_+,J_-\in \Bbb{CP}^1_h$，此时 \cite{ref67} 的公式略有简化，因为 $B$-场可以取为零（之后再通过 $B$-场变换恢复）。令 $\omega_\pm$ 为两个辛结构 $\omega_\pm = gJ_\pm$。于是按 \cite{ref67} 的式 6.3，由 $J_+$、$J_-$ 与超凯勒度量 $g$ 这组数据确定的两个广义复结构之一是 % 式 (5.27)
\setcounter{equation}{26}
\begin{equation}\label{eq:ugg}
{\cal J}={1\over 2}\left(\begin{matrix}J_++J_- &
-(\omega_+^{-1}-\omega_-^{-1})\\ \omega_+-\omega_- &
-(J_+^t+J_-^t)\\\end{matrix}\right).
\end{equation}
第二个与我们无关，可通过把 $J_-$ 与 $\omega_-$ 反号得到。

利用 \eqref{eq:wyrre}\ 以及 $(w_+,w_-)=(-t,t^{-1})$，我们有 % 式 (5.28)
\setcounter{equation}{27}
\begin{equation}\label{eq:ytro}
\begin{aligned} J_+ & = {1-\bar t t\over 1+\bar t t}I-{i(t-\bar t)\over
1+\bar t t}J-{t+\bar t\over 1+\bar t t}K  \\ J_- & = -{1-\bar
tt\over 1+\bar t t}I-{i(t-\bar t)\over 1+\bar t t}J+{t+\bar t\over
1+\bar t t}K  .\\\end{aligned}
\end{equation}
相应的辛结构为
 % 式 (5.29)
\setcounter{equation}{28}
\begin{equation}\label{eq:bytro}
\begin{aligned} \omega_+ &={\rm Im}\,\tau\left( {1-\bar t t\over 1+\bar
t t}\omega_I-{i(t-\bar t)\over 1+\bar t t}\omega_J-{t+\bar t\over
1+\bar t t}\omega_K\right) \\ \omega_- & = {\rm
Im}\,\tau\left(-{1-\bar tt\over 1+\bar t t}\omega_I-{i(t-\bar
t)\over 1+\bar t t}\omega_J+{t+\bar t\over 1+\bar t
t}\omega_K\right) .\\\end{aligned}
\end{equation}
（${\rm Im}\,\tau$ 这一因子来自 \eqref{eq:inog}。）于是我们算出由这些数据确定的广义复结构为 % 式 (5.30)
\setcounter{equation}{29}
\begin{equation}\label{eq:islo}
{\cal I}_t={1\over
1+\bar t t}\left(\begin{matrix}-i(t-\bar t)J & -({\rm
Im}\,\tau)^{-1}((1-\bar t t)\omega_I^{-1}-(t+\bar
t)\omega_K^{-1})\\ {\rm Im}\,\tau ((1-\bar t t)\omega_I-(t+\bar
t)\omega_K)& i(t-\bar t)J^t\\\end{matrix}\right).
\end{equation}
（要得到这个公式，有用的知识是：在超凯勒流形上，若 $a^2+b^2+c^2=1$，则 $(a\omega_I+b\omega_J+c\omega_K)^{-1}=a\omega_I^{-1}+b\omega_J^{-1}+c\omega_K^{-1}$。）

第一个重要的观察是：$I$、$K$ 与 $\omega_J$ 都消失了。${\cal I}_t$ 只依赖于 $J$、$\omega_K$ 与 $\omega_I$，或等价地只依赖于 $J$ 与全纯二形式 $\Omega_J=\omega_K+i\omega_I$（及其复共轭）。但正如我们在第~\ref{sec:mhhyper}~节所指出的，这些数据与 $C$ 的度量无关。这就从二维的观点解释了：为什么由四维维数约化得到的这族拓扑场论不依赖于度量。

\begin{yizhu}[易错点]
按：${\cal I}_t$ 中 $I$ 与 $\omega_J$ 之所以「消失」，是因为 $J_\pm$、$\omega_\pm$ 的组合把这两项相消掉了；剩下的数据恰好是 $J$ 与 $\Omega_J=\omega_K+i\omega_I$。因此这族拓扑场论与 $C$ 的度量的无关性，根源在于超凯勒结构及其全纯二形式 $\Omega_J$ 都不依赖 $C$ 的度量，而不在于某个近似或特殊取值。
\end{yizhu}

在任意超凯勒流形 $X$ 上也可构造类似的广义复结构族，\cite{ref67} 第 4.6 节对此有描述。那里（命题 4.34）解释说，这样一族中的一般元素是某个由辛结构导出的广义复结构的 $B$-场变换。的确，稍作计算就得到
% 式 (5.31)
\setcounter{equation}{30}
\begin{equation}\label{eq:huffy}
{\cal I}_t={\cal M}(B_0)\left(\begin{matrix}0 &
-\omega_0^{-1}\\ \omega_0 & 0\\\end{matrix}\right){\cal M}(B_0)^{-1},
\end{equation}

其中 % 式 (5.32)
\setcounter{equation}{31}
\begin{equation}\label{eq:nutty}
\begin{aligned}\omega_0 & = {\rm Im}\,\tau{1-\bar
t{}^2t^2\over (1+t^2)(1+\bar t{}^2)}\left(\omega_I-{t+\bar t\over
1-\bar t t}\omega_K\right)\\
                      B_0 &= -{\rm Im}\,\tau\left.{i(t^2-\bar t{}^2)\over(1+t^2)(1+\bar
                      t{}^2)}\right.
                      \left(\omega_I+{1-\bar t t\over t+\bar t}\omega_K\right).\\\end{aligned}
\end{equation}

因此，对于 $t\neq \pm i$，由广义复结构 ${\cal I}_t$ 导出的 TQFT 等价于具有辛结构 $\omega_0$ 与 $B$-场 $B_0$ 的 $A$-模型。这一事实将使我们能从二维的观点理解第~\ref{sec:canonpar}~节引入的典范参数 $\Psi$。

一般地，具有辛形式 $\omega_0$ 与 $B$-场 $B_0$ 的 $A$-模型只依赖于上同调类 $[B_0+i\omega_0]$。（在 Kähler 流形上，这称为复化凯勒类。）在本问题中，由于 $\omega_K$ 的上同调类为零（正如我们在第~\ref{sec:mhhyper}~节所解释的），我们有 % 式 (5.33)
\setcounter{equation}{32}
\begin{equation}\label{eq:gutty}
[B_0+i\omega_0]=-i\,{\rm
Im}\,\tau\left(t-t^{-1}\over t+t^{-1}\right) [\omega_I].
\end{equation}
因此，当四维 $\theta$ 角为零时（我们通过从 ${\cal I}_t$ 而非其 $B$-场变换出发来取它为零），模型只通过 \eqref{eq:gutty} 中出现的那种组合依赖于 ${\Im}\tau$ 与 $t$。

如我们在第~\ref{sec:mhhyper}~节末尾所解释的，四维 $\theta$ 角会给二维 $B$-场带来额外贡献 $B'=-(\theta/2\pi)\omega_I=-({\rm Re}\,\tau)\,\omega_I$。这可以简单地通过用 ${\cal M}(B')$ 作共轭而纳入广义复结构。所得的模型依赖于上同调类 $[B_0+B'+i\omega_0]$，它等于 $-[\omega_I]$ 乘以
% [WARN] 公式 \eqn\eqref{eq:loopy} 无编号或重复
\begin{equation*}
\tag*{(5.34)}
\Psi={\rm Re}\,\tau +i\,{\rm Im}\,\tau
\left(t-t^{-1}\over t+t^{-1}\right).
\end{equation*}
这就从二维的观点解释了模型为何只通过组合 $\Psi$ 依赖于 $t$ 与 $\tau$。

\begin{yizhu}[含糊点拨]
按：$\Psi$ 的引入逻辑是：$A$-模型只依赖复化凯勒类 $[B_0+i\omega_0]$，而在计入 $\theta$ 角贡献的 $B'=-({\rm Re}\,\tau)\omega_I$ 之后，这个类正比于 $-[\omega_I]$ 乘以组合式 (5.34)。所以「只依赖 $\Psi$」并非新的动力学输入，而是四维 $\theta$ 角与二维 $B$-场变换的记账结果；$t=\pm i$ 处 $\Psi=\infty$，此推导失效，须单独处理。
\end{yizhu}

\par\medskip{\it 几处零碎的收尾}

最后，我们来收拾几处未了的细节。

首先，若 $t$ 为实数，则由复结构 $J_+(t),J_-(t)$ 确定的 TQFT 本身就是 $A$-模型，而上面证明 $\Psi$ 是唯一相关参数的论证其实并不需要广义复几何。我们本可以基于实情形并利用对 $t$ 的全纯性来推理。我们选择不这样做，是因为我们认为广义复几何是理解这一问题的自然框架，尽管目前对它的理解尚存一些技术上的缺口。

其次，在上面的分析中我们用了「$A$-模型只依赖于 $B+i\omega$ 的上同调类」这一事实。其证明是把作用量写成 $\{Q,V\}+\int_\Sigma \Phi^*(\omega-iB)$，其中 $\Phi:\Sigma\to X$ 是 sigma 模型的映射。

在上述推导中，$B$-场并不只是 $(1,1)$ 型的（在使辛形式成为 Kähler 形式的那个复结构中）；它还有 $(2,0)\oplus (0,2)$ 型的分量。事实上，$A$-模型对 $B$-场的所有分量都敏感，包括 $(2,0)\oplus (0,2)$ 型的那部分，但这一点可能需要澄清。

若 $\Phi:\Sigma\to X$ 是全纯映射，而 $B$ 是 $X$ 上 $(2,0)\oplus (0,2)$ 型的形式，则 $\Phi^*(B)=0$。于是，如果人们把 $A$-模型纯粹解释为通过对全纯映射求和来计算关联函数的机制，$B$ 的 $(2,0)\oplus (0,2)$ 部分就显得无关紧要。但若考虑膜（为了理解几何 Langlands 纲领，我们肯定会这样做），人们立刻会看到 $B$ 的 Hodge 分解的所有部分都相关。带有曲率 $F$ 的 Chan-Paton 线丛 ${\cal L}$ 的 Lagrangian 子流形 $N\subset X$ 给出一个 $A$-膜，当且仅当 $F+B|_N=0$。这个条件显然对 $B$ 的 $(2,0)\oplus (0,2)$ 部分敏感。

最后，我们关于 $\Psi$ 是唯一相关参数的分析在 $t=\pm i$ 时其实并不成立，因为公式在该处有极点。在这些 $t$ 值下，模型实际上不是某个 $A$-模型的 $B$-场变换，而是复结构 $\pm J$ 中的 $B$-模型。为补全这些 $t$ 值下的分析，我们将直接用 $B$-模型来论证。

在 $t=\pm i$ 时 $\Psi=\infty$，与 $\tau$ 无关；因此为完成「$\Psi$ 是唯一相关参数」的证明，我们必须证明 $\tau$ 在 $t=\pm i$ 时完全无关。${\Im}\,\tau$ 控制 $\MH$ 的 Kähler 类，而后者在 $B$-模型中当然无关。改变 ${\rm Re}\,\tau$ 会给 $B$-场加上 $\omega_I$ 的一个倍数。为说明该项在复结构 $J$ 的 $B$-模型中无关，我们写 $\omega_I=-i\Omega_J+i\omega_K$。来自 $\omega_K$ 的贡献无关，因为形式 $\omega_K$ 是恰当的。来自 $\Omega_J$ 的贡献无关，因为 $\Omega_J$ 是 $(2,0)$ 型形式。而在 $B$-模型中，若 $B$ 是 $(1,1)$ 型或 $(2,0)$ 型，$B$-场贡献 $-i\int_\Sigma\Phi^*(B)$ 可以写成 $\{Q,\dots\}$。（相反，$B$-场的 $(0,2)$ 分量确实会影响 $B$-膜的范畴 \cite{ref71}。）当然，在复结构 $-J$ 中，我们从 $\omega_I=i\bar\Omega_J-i\omega_K$ 出发作同样的论证。

\begin{yizhu}[易错点]
按：$t=\pm i$ 是两个特殊值：此时模型是复结构 $\pm J$ 中的 $B$-模型，而非某个 $A$-模型的 $B$-场变换，因此 $\Psi$ 的推导在极点处不适用。另一个容易忽略的细节是 $B$-场的 $(0,2)$ 分量在 $B$-模型中会影响 $B$-膜的范畴，而 $(1,1)$、$(2,0)$ 分量不进入范畴。
\end{yizhu}

\subsection{一些特例}

\label{sec:somespec}

在本文余下部分（除第~\ref{sec:genotics}~节外），我们主要关注特殊值 $\Psi=0$ 与 $\Psi=\infty$，以及对偶 % 式 (5.35)
\setcounter{equation}{34}
\begin{equation}\label{eq:rtu}
S=\left(\begin{matrix}0 & 1/\sqrt {n_{\frak
g}}\\ -\sqrt {n_{\frak g}}& 0\\\end{matrix}\right)
\end{equation}
它把 $\Psi\to -1/n_{\frak g}\Psi$，从而交换 $\Psi=0$ 与 $\Psi=\infty$。我们回忆，对单连通型 $G$ 有 $n_{\frak g}=1$。当我们说 $S$-对偶时，对于指的是由 $S$ 与 % 式 (5.36)
\setcounter{equation}{35}
\begin{equation}\label{eq:ntu}
T=\left(\begin{matrix}1 & 1\\ 0& 1\\\end{matrix}\right),
\end{equation}
生成的整个对偶群 $\Gamma$、还是仅指由 $S$ 生成的子群，我们会有少许不精确。事实上，完整的故事还需要把 $T$ 包括进来，但在几何 Langlands 纲领中 $S$ 才是主要角色。

为得到 $\Psi=\infty$，我们取 $t=i$，从而 $(w_+,w_-)=(-i,-i)$。所得的模型是复结构 $J$ 中的 $B$-模型。这一陈述不受 $\tau$ 选取的影响，但取 $\tau$ 为虚数较为方便。

施加变换 $S$，我们发现 $\tau$ 仍为虚数，而 $t$ 被映为 $t'=-t({\tau}/|\tau|)=-it=1$。这导致 $(w_+,w_-)=(-1,1)$，从而是复结构 $K$ 中的 $A$-模型。因此 $S$-对偶蕴含：复结构 $J$ 中 $\MH(G,C)$ 的 $B$-模型等价于复结构 $K$ 中 $\MH(\LG,C)$ 的 $A$-模型。若把 $S$ 与把复结构 $K$ 转回复结构 $J$ 的变换 ${\EUI}:\varphi\to i\varphi$ 结合，我们得到把复结构 $J$ 中的 $B$-模型映到复结构 $J$ 中的 $A$-模型的镜像对称。于是复流形 $\MH(G,C)$ 与 $\MH(\LG,C)$ 在复结构 $J$ 中构成一对镜像。

然而，我们这里用到的自同构 ${\EUI}$ 并不十分自然，因为它把我们移出了由 $\Bbb{CP}^1_g$ 参数化的四维 TQFT 族。它必须把我们移出这个族，因为复结构 $J$ 中的 $A$-模型并不是这个族的成员！

\begin{yizhu}[物理图像]
按：在二维图像里，$S$-对偶把 $J$ 中的 $B$-模型送到 $K$ 中的 $A$-模型（见下表），再复合自同构 ${\EUI}$（它把 $K$ 转回 $J$，代价是离开 $\Bbb{CP}^1_g$ 族）就得到 $J$ 中的镜像对称。于是 $S$-对偶的一部分可视为镜像对称的「一半」，这也是 $S$ 在几何 Langlands 纲领中扮演核心角色的根源。
\end{yizhu}

此外，我们在讨论 \eqref{eq:yfog}\ 时已经注意到，$S$（以及整个对偶群）保持复结构 $I$ 中的 $B$-模型不变，同样也保持该复结构中的 $A$-模型不变。最后，既然 $S$ 把复结构 $J$ 中的 $B$-模型与复结构 $K$ 中的 $A$-模型互换，我们可以猜测它同样把复结构 $J$ 中的 $A$-模型与复结构 $K$ 中的 $B$-模型互换。事实上，自同构 ${\EUI}$ 把这两个陈述联系起来，但我们不这样论证，而改用式 \eqref{eq:yfog}。复结构 $J$ 中的 $A$-模型对应于 $(w_+,w_-)=(-i,i)$。它被 $S$ 映到 $(w_+,w_-)=(1,1)$，而后者确实对应于复结构 $K$ 中的 $B$-模型。我们把上述陈述汇总成一张表。

\par\medskip

\begin{table}[htbp]
\centering
\begin{tabular}{ll}
\toprule
模型 & $S$-对偶 \\
\midrule
$I_B$ & $I_B$ \\
$I_A$ & $I_A$ \\
$J_B$ & $K_A$ \\
$J_A$ & $K_B$ \\
$K_B$ & $J_A$ \\
$K_A$ & $J_B$ \\
\bottomrule
\end{tabular}
\caption{一个模型及其 $S$-对偶。例如 $I_B$ 表示复结构 $I$ 中的 $B$-模型，它是自身的 $S$-对偶。（原文 Table~2）}
\label{tab:sdual-models}
\end{table}

\begin{yizhu}[易错点]
按：表中下标 $B$/$A$ 指模型类型，与复结构记号 $I,J,K$ 搭配：$J_B$ 是「复结构 $J$ 中的 $B$-模型」。$I_B$、$I_A$ 是自身的 $S$-对偶，这来自 \eqref{eq:yfog}（$S$ 保持 $I$ 中的 $B$-模型与 $A$-模型不变）；同时不要把复结构 $I$ 与广义复结构 ${\cal I}$ 混淆。
\end{yizhu}
